\documentclass[11pt,a4paper]{article}
\usepackage[margin=1.8cm]{geometry}
\usepackage{amsmath,amssymb,tabularx,array,needspace,fontspec,enumitem,cancel}
\newcommand{\tb}{\cancel{||||}\,}
\setmainfont{FreeSerif}
\newfontfamily\urdufont[Script=Arabic]{FreeSerif}
\TeXXeTstate=1
\newcommand{\urdu}[1]{{\urdufont\beginR #1\endR}}
\setlength{\parindent}{0pt}
\setlength{\parskip}{0pt}
\renewcommand{\arraystretch}{1.15}
\pagestyle{plain}

% \mcq{number}{question}{A}{B}{C}{D}
\newcommand{\mcq}[6]{%
  \Needspace{4\baselineskip}%
  \par\medskip\noindent\textbf{#1.}\ #2\par\nopagebreak
  \noindent\begin{tabularx}{\linewidth}{@{\hspace{1.5em}}*{4}{>{\raggedright\arraybackslash}X}@{}}
    (A)~#3 & (B)~#4 & (C)~#5 & (D)~#6
  \end{tabularx}\par}

\begin{document}

\begin{center}
  {\Huge\bfseries The Hope Science Academy}\\[3pt]
  {\small Excellence is our standard not goal}\\[2pt]
  {\small 247-E-1, (Opposite LDA School) Johar Town, Lahore. PH: 0300-8194789, 0322-4157001}
\end{center}

\medskip
\noindent
\begin{tabularx}{\linewidth}{|X|X|X|X|}
\hline
\textbf{Student Name:} & \textbf{Roll Num:} & \textbf{Class:} 12th & \textbf{Paper Code:} \\ \hline
\textbf{Subject Name:} Principles of Accounting & \textbf{Time Allowed:} & \textbf{Total Marks:} 525 & \textbf{Exam Date:} 02-Nov-2026 \\ \hline
\multicolumn{4}{|l|}{\textbf{Exam Syllabus:} FULL BOOK} \\ \hline
\end{tabularx}

\bigskip
\noindent\textbf{Q1. Choose the correct answer.}\hfill\textbf{1$\times$525=525}

\mcq{1}{Double entry system means:}{The recording of same things twice}{Double accounting}{The two sides of every transaction are recorded}{Entry at two dates}
\mcq{2}{Under single entry system sometimes:}{Double aspects of a transaction are recorded}{Single aspect of a transaction is recorded}{No aspects of a transaction is recorded}{All of the above}
\mcq{3}{Single entry system is generally adopted by:}{Small business concerns}{Large business concerns}{Medium business concerns}{All of the above}
\mcq{4}{Single entry system contains:}{A daybook or general journal}{A cashbook}{Ledger accounts for individual customers and creditors}{All of above}
\mcq{5}{In single entry system:}{Personal accounts are maintained}{Real accounts are maintained}{Nominal accounts are maintained}{None of the above}
\mcq{6}{In single entry system ledger contains the accounts of:}{Debtors, creditors and cash}{Wages and carriage}{Salaries and insurance}{Assets and liabilities}
\mcq{7}{In single entry system nominal accounts and some assets accounts balance are not available so it is not possible to prepare:}{Debtors account}{Creditors account}{Trial balance}{Statement of affairs}
\mcq{8}{In single entry system the final accounts are prepared by means of statement of affairs:}{Correctly}{EffectivelyEffectively}{Wrongly}{Approximately}
\mcq{9}{In single entry system real accounts are not maintained so the preparation of:}{Opening statement of affairs is not possible}{Closing statement of affairs is not possible}{Balance sheet is not possible}{“a” and “b”}
\mcq{10}{How many methods are available to calculate the profit or loss under single entry system:}{One}{Two}{Three}{Four}
\mcq{11}{Net worth is equal to:}{Assets – Liabilities}{Liabilities – Assets}{Liabilities + Capital}{Capital + Assets}
\mcq{12}{To calculate the value of capital both at the beginning and at the end of year in single entry system we prepare:}{Balance sheet}{Trading account}{Trial balance}{Statement of affairs}
\mcq{13}{Opening capital is calculated by taking the difference between:}{Opening assets and closing assets}{Opening assets and closing liabilities}{Opening assets and Opening liabilities}{Closing assets and closing liabilities}
\mcq{14}{Closing capital in single entry system is calculated by preparing:}{Opening statement of affairs}{Closing statement of affairs}{Opening balance sheet}{Closing balance sheet}
\mcq{15}{Due to fresh capital introduced during the year, the closing capital will:}{Increase}{Decrease}{Constant}{Multiply}
\mcq{16}{Due to drawings made during the year, the closing capital will:}{Increase}{Decrease}{Constant}{Multiply}
\mcq{17}{To calculate the true net profit or net loss in single entry system it is necessary that the amount of drawings should be added in:}{Fresh capital}{Opening capital}{Closing capital}{Adjusted closing capital}
\mcq{18}{To calculate the true net profit or net loss in single entry system it is necessary that the amount of fresh capital should be subtracted from:}{Adjusted closing capital}{Closing capital}{Cash in hand}{Cash at bank}
\mcq{19}{Single entry system is situated where:}{Small amount of cash transactions are more}{Large amount of cash transactions are more}{Small amount of credit transactions are more}{large amount of credit transactions are more}
\mcq{20}{According to companies ordinance, single entry system cannot be adopted by:}{Sole proprietorship}{Partnership}{“a” and “b”}{Joint stock companies}
\mcq{21}{Both aspects of a transaction must be recorded in:}{Single entry system}{Barter system}{Double entry system}{Islamic system}
\mcq{22}{In increased net worth method, profit or loss is calculated by preparing:}{Income and expenditure account}{Profit and loss account}{Statement of profit or loss}{Balance sheet}
\mcq{23}{We cannot prepare the trial balance and precise balance sheet under:}{Single entry system}{Double entry system}{“a” and “b”}{Barter system}
\mcq{24}{In single entry system, statement of assets, liabilities and capital is called:}{Income statement}{Earning statement}{Statement of profit or loss}{Statement of affairs}
\mcq{25}{The format sketch of statement of affairs is similar to:}{Trial balance}{Profit \& loss account}{Balance sheet}{Income \& Expenditure account}
\mcq{26}{It is prepared to calculate the net profit or net loss under single entry system:}{Opening statement of affairs}{Statement of profit or loss}{Receipts and payment account}{Balance sheet}
\mcq{27}{If the adjusted capital at the end is more than the capital in the beginning then the difference will be:}{Gross profit}{Gross loss}{Net loss}{Net profit}
\mcq{28}{If the capital at the end Rs.30000; drawings Rs.1250; capital in the beginning Rs.15000; then the amount of net profit will be:}{Rs.43750}{Rs.16250}{Rs.13750}{Rs.43500}
\mcq{29}{If the capital at the end Rs.45000; fresh capital introduced during the year Rs.15000; capital in the beginning Rs.32500; then the amount of loss will be:}{Rs.2500}{Rs.27500}{Rs.62500}{Rs.45000}
\mcq{30}{If the capital at the end Rs.125000 ; drawings Rs.12000; capital in the beginning Rs.60000; profit during the year Rs.27000; then the amount of fresh capital introduced will be:}{Rs.87000}{Rs.50000}{Rs.77000}{Rs.170000}
\mcq{31}{If the capital at the end Rs.55000; capital in the beginning Rs.56375; loss during the year Rs.1375; then the amount of drawings will be:}{Rs.110000}{Rs.2750}{Nill}{Rs.112750}
\mcq{32}{If the capital in the beginning Rs.65650; fresh capital introduced during the year Rs.16250; loss during the year Rs.11375; then the capital at the end will be:}{Rs.70525}{Rs.93275}{Rs.38025}{Rs.48075}
\mcq{33}{If the capital at the end Rs.82500; drawings Rs.12925; profit during the year Rs.16500; then the capital in the beginning will be:}{Rs.111925}{Rs.86075}{Rs.53075}{Rs.78925}
\mcq{34}{If the capital in the beginning 01-01-2016 Rs.80000; capital introduced on 01-05-2004 Rs.20000; rate of interest 6\% P.a., then the interest on capital at 31-12-2004 will be:}{5600Rs.5600}{4800Rs.4800}{6000Rs.6000}{4000Rs.4000}
\mcq{35}{The opening balance of debtors can be ascertained by preparing:}{Total creditors’ account}{Cash account}{Balance sheet}{Total debtors’ account}
\mcq{36}{The amount of credit purchases can be ascertained by preparing:}{Total debtors’ account}{Total creditors’ account}{Sales account}{Income and expenditure account}
\mcq{37}{The amount of cash purchases and cash sales can be ascertained by preparing:}{Total debtors’ account}{Total creditors’ account}{Profit \& loss account}{None of the above}
\mcq{38}{Cash paid to creditors can be ascertained by preparing:}{Total debtors’ account}{Total creditors’ account}{Balance sheet}{None of the above}
\mcq{39}{The closing balance of creditors can be ascertained by preparing the:}{Bills receivable account}{Bill payable account}{Total debtors’ account}{Total creditors’ account}
\mcq{40}{The opening and closing balance of bills receivable can be calculated by preparing the:}{Sales account}{Total debtors’ account}{Total creditors’ account}{Bills receivable account}
\mcq{41}{If sales Rs.90000; Purchase Rs.65625; closing stock Rs.11250; rate of gross profit 20\% on sales, then the stock in the beginning will be:}{Rs.17625}{Rs.18625}{Rs.16625}{Rs.15625}
\mcq{42}{If sales Rs.45000; Purchase Rs.65000; opening stock Rs.3750; gross profit on sales 20\%, then the closing stock will be:}{Rs.28250}{Rs.28260}{Rs.28225}{Rs.28275}
\mcq{43}{If opening stock Rs.15000; sales Rs.180000; purchase Rs.140000; gross profit on cost 20\%, then the closing stock will be:}{Rs.10000}{Rs.150000}{Rs.5000}{Nill}
\mcq{44}{If cost of goods sold Rs.350000; gross profit on cost 15\%, then the value of sales will be:}{Rs.42500}{Rs.402500}{Rs.42250}{Rs.40500}
\mcq{45}{If cost of goods sold Rs.1,400,000; gross profit on cost 205\%, then the value of sales will be:}{Rs.1,166,665}{Rs.1,166,667}{Rs.1,166,660}{Rs.1,166,000}
\mcq{46}{If opening balance of debtors Rs.26500; closing balance of debtors Rs.44000; credit sales Rs.44000; bad debts Rs.2000, then the value of cash received from debtor will be:}{Rs.24000}{Rs.25000}{Rs.24500}{Rs.25500}
\mcq{47}{If cash paid to creditors Rs.8750; closing balance of creditors Rs.8535; bills payable Rs.2500, then the value of credit purchase will be:}{Rs.14785}{Rs.19785}{Rs.2715}{Rs.2285}
\mcq{48}{If cash paid to creditors Rs.33075; returns outwards Rs.225; credit purchase Rs.67500; cash purchases Rs.25000, then the value of closing creditors will be:}{34300Rs.34300}{34500Rs.34500}{34200Rs.34200}{34000Rs.34000}
\mcq{49}{Net profit making organizations:}{Buy goods}{Manufacturing goods}{Sell goods}{None of these}
\mcq{50}{Example of non-profit making organization:}{Fan factory}{Sugar industry}{Hospital}{Shoe factory}
\mcq{51}{Non-profit making organizations are established for:}{Profit}{Charitable or religious purpose}{To manufacture goods}{To help the rich people}
\mcq{52}{Non-profit making organization is also known as:}{Non-profit seeking concern}{Non-trading concern}{Trading concern}{“a” and “b”}
\mcq{53}{The method for preparing the final accounts in trading concerns and non- profit making organizations are:}{Same}{Different}{Easy}{Difficult}
\mcq{54}{A person who is elected to run the club:}{Treasure}{Secretary}{Chairperson}{President}
\mcq{55}{The basic club record or book of original entry in non-profit making organization:}{Purchase book}{Sales book}{Cash book}{“a” and “b”}
\mcq{56}{Receipts and payments account is a summary of:}{Cash book}{Purchase book}{Sales book}{Purchase returns book}
\mcq{57}{Receipts and payments account records the transactions of:}{Revenue nature}{Capital nature}{Revenue and capital nature}{None of the above}
\mcq{58}{Receipts and payments account shows:}{Opening and closing balances of cash}{Assets and liabilities}{Liabilities and capital}{Profit and loss}
\mcq{59}{Receipts and payments account shows:}{Debit and credit balances of ledgers}{Income and expenditures}{Cash receipts and payments}{Assets and liabilities}
\mcq{60}{A receipts and payments account is similar to:}{An income and expenditure account}{A statement of affairs}{A cash or a bank account}{A profit or loss account}
\mcq{61}{Receipts and payments account records transactions relating to:}{Past year}{Future year}{Present year}{All of the above}
\mcq{62}{In receipts and payments account:}{No narration is written}{No ledger reference is recorded}{Daily cash balance cannot be ascertained}{All of the above}
\mcq{63}{In cash book:}{Each entry is followed by narration}{Ledger reference is recorded}{Each transaction is recorded in chronological}{All of the above}
\mcq{64}{In receipts and payment account:}{Credit transaction are not recorded}{Non-cash items are recorded}{“a” and “b”}{Cash transaction are not recorded}
\mcq{65}{In non-profit making organization final accounts contains:}{Income and expenditure account and balance sheet}{Receipts and payments account and balance sheet}{Receipts and payments account and statement of affairs}{Profit and loss account and balance sheet}
\mcq{66}{It is prepared by non-profit making organization to determine surplus or deficit:}{Trading account}{Profit and loss account}{Receipts and payment account}{Income and expenditure account}
\mcq{67}{Income and expenditure account records the transactions of:}{Revenue nature}{Capital nature}{Revenue and capital nature}{None of these}
\mcq{68}{Non-profit seeking organizations prepare:}{Manufacturing account}{Trading account}{Income and expenditure account}{Profit and loss account}
\mcq{69}{Non-profit seeking organization income and expenditure account performs the same function as a profit seeking organization’s:}{Trading and profit and loss account}{Balance sheet}{Bank account}{Petty cash book}
\mcq{70}{Income and expenditure account shows:}{Cash in hand}{Surplus or deficit}{Cash at bank}{Capital account}
\mcq{71}{All transactions relating to non-profit seeking concerns are recorded in the books of account strictly according to:}{Single entry system}{Double entry system}{Trading entry system}{“a” and “b”}
\mcq{72}{In income and expenditure account:}{Expenses are recorded on debit side}{Revenues are recorded on credit side}{Does not start with any opening balance}{All of the above}
\mcq{73}{Capital fund is also called:}{General fund}{Surplus fund}{Accumulated fund}{All of these}
\mcq{74}{The amount received by the concern as per the will of the donor is known as:}{Legacy}{Donation}{Subscription}{Expense}
\mcq{75}{Subscription received relating to current year is:}{An income}{An expense}{An asset}{A liability}
\mcq{76}{Subscription received in advance is:}{An income}{An expense}{A liability}{An asset}
\mcq{77}{Subscription received during the current year is:}{An income}{An expense}{An asset}{a \& c}
\mcq{78}{Amount received from any source by way of gift is described as:}{Legacy}{Subscription}{Donation}{Life membership fee}
\mcq{79}{Lump sum amount received from members which enable them to become the members for whole of the life:}{Life member ship fee}{Donation}{Subscription}{Legacy}
\mcq{80}{If the amount of life membership fees is small and it is due to regular flow of members then it:}{Should be debited to income and expenditure account}{Should be credited to income and expenditure account}{Should be debited to a special fund}{Should be credited to a special fund}
\mcq{81}{The amount paid by a new partner at the time of admission:}{Admission fee}{Entrance fee}{Life membership fee}{a \& b}
\mcq{82}{The amount paid to persons who are invited to deliver lectures in a club is known as:}{Salary}{Wages}{Honorarium}{Income}
\mcq{83}{Amount of honorarium should be:}{Debited to income and expenditure account}{Credited to income and expenditure account}{Debited to special fund}{Credited to special fund}
\mcq{84}{Sales of newspaper should be:}{Debited to income and expenditure account}{Credited to income and expenditure account}{Debited to asset account}{Credited to special fund}
\mcq{85}{Sales of sports material be treated as:}{Income}{Expense}{Asset}{Liability}
\mcq{86}{If there is tournament fund, then tournament expense and incomes are transferred to:}{Income and expenditure account}{Asset side of balance sheet}{The liability side of the balance sheet}{Both income and expenditure account and balance sheet}
\mcq{87}{Life membership fee Rs.10000, 90\% capitalized should be recorded in:}{Receipts and payments account}{Income and expenditure account}{Balance sheet}{All of these}
\mcq{88}{Donations Rs.50,000, 100\% capitalized should be recorded in:}{Receipts and payments account}{Income and expenditure account}{a \& b}{Receipts and payments account and balance sheet}
\mcq{89}{Entrance fee Rs.8,000 which is 90\% capitalized, the amount should be credited to income and expenditure account is:}{Rs.800}{Rs.7,600}{Rs.7,200}{Rs.8,000}
\mcq{90}{Donations Rs.9000, 4/5th capitalized, the amount should be credited to income and expenditure account is:}{1800Rs.1,800}{7200Rs.7,200}{9000Rs.9,000}{900Rs.900}
\mcq{91}{If salaries Rs.300 paid during the year and Rs.25 is outstanding, then the amount debited to income and expenditure account will be:}{300Rs.300}{325Rs.325}{275Rs.275}{350Rs.350}
\mcq{92}{If Rs.700 paid for interest during the year and amount debited to income and expenditure account is Rs.780, then the amount of outstanding interest will be:}{700Rs.700}{780Rs.780}{80Rs.80}{740Rs.740}
\mcq{93}{If the amount debited to income and expenditure account is Rs.450 for printing and stationery and outstanding is Rs.75, then the amount paid during the year will be:}{375Rs.375}{325Rs.325}{525Rs.525}{630Rs.530}
\mcq{94}{If salary paid during the year is Rs.1200 and prepaid salaries is Rs.100, then the amount debited to income and expenditure account will be:}{1200Rs.1,200}{1100Rs.1,100}{1300Rs.1,300}{1500Rs.1,500}
\mcq{95}{If travelling expenses paid during the year Rs.1200 and amount debited to income and expenditure account is Rs.1,170 then the amount for prepaid travelling expenses will be:}{60Rs.60}{1130Rs.1,130}{1230Rs.1,230}{30Rs.30}
\mcq{96}{If rent received during the year is Rs.5000 and the amount of pre-receive rent is Rs.600 then the amount credited to income and expenditure account will be:}{4400Rs.4,400}{5600Rs.5,600}{5000Rs.5,000}{6200Rs.6,200}
\mcq{97}{If the amount of accrued commission is Rs.400 and the amount of commission received is Rs.2,800 then the amount credited to income and expenditure account will be:}{2800Rs.2,800}{400Rs.400}{2400Rs.2,400}{3200Rs.3,200}
\mcq{98}{If salaries paid during the year is Rs.1,575 and amount debited to income and expenditure account is Rs.1,500 then the amount prepaid salaries will be:}{75Rs.75}{1500Rs.1,500}{1575Rs.1,575}{1425Rs.1,425}
\mcq{99}{If the amount of prepaid rent is Rs.150 and amount debited to income and expenditure account is Rs.3250 then rent paid during the year will be:}{3400Rs.3,400}{1300Rs.1,300}{3500Rs.3,500}{3000Rs.3,000}
\mcq{100}{Receipt and payment account is a:}{Nominal account}{Personal account}{Real account}{Impersonal account}
\mcq{101}{Donation received for special purpose:}{Should be credited to income and expenditure account}{Should be debited to income and expenditure account}{Should be credited to a separate account and shown in the balance sheet}{Should be debited as separate account and shown in the balance sheet}
\mcq{102}{Payment of Rs. 2500 on the purchase of billiard table should be entered in:}{Receipts and payment account}{Receipts and payment account and income and expenditure account}{Receipts and payment account and balance sheet}{Income and expenditure account and balance sheet}
\mcq{103}{Depreciation written off on the furniture should be entered in the:}{Receipts and payments account}{Receipts and payments account and income and expenditure account}{Receipts and payments account and balance sheet}{Income and expenditure account and balance sheet}
\mcq{104}{Entrance fee, whole amount credited to income and expenditure account should be recorded in:}{Income and expenditure account}{Receipts and payment account and income and expenditure account}{Receipts and payments account and balance sheet}{All of these}
\mcq{105}{In accountancy ‘consignment’ signifies:}{Goods forwarded from one place to another}{Goods forwarded by a person to another}{Goods dispatched by its owner to his agent}{Goods dispatched by its owner to his agent for the purpose of sale}
\mcq{106}{The parties involved in consignment are:}{Two}{Three}{Four}{Five}
\mcq{107}{The manufacturer or wholesaler who sends his goods for sale purpose is known as:}{Consignee}{Consignor}{Agent}{Drawer}
\mcq{108}{Consignee is:}{Principal}{Owner}{Debtor of the consignee}{None of these}
\mcq{109}{The person to whom the goods are sent is known as:}{Consignor}{Consignee}{Drawer}{Drawee}
\mcq{110}{Consignee is:}{Principal}{Agent}{Creditor of the consigner}{None of the above}
\mcq{111}{Consignee sells the goods on behalf of:}{Drawer}{Drawee}{Consignor}{Agent}
\mcq{112}{The relationship between consignor and consignee is that of:}{Principal and agent}{Debtor and creditor}{Buyer and seller}{None of these}
\mcq{113}{The act of sending goods by owner to his agent for sale is called:}{Consignment}{Sales}{Purchases}{Partnership}
\mcq{114}{When goods are dispatched by the consigner to the consignee, it will be from consignor’s view point:}{Consignment inward}{Consignment outward}{Consignment in transit}{All of these}
\mcq{115}{When goods are dispatched by the consignor to the consignee it will be from consignee’s view point:}{Consignment outward}{Consignment in past}{Consignment in future}{Consignment inward}
\mcq{116}{The register which is maintained by consignor to keep full record of the goods consigned is called:}{Consignment outward book}{Consignment inward book}{Consignment production book}{Consignment in purchases book}
\mcq{117}{Consignment outward book treatment is similar to:}{Purchase book}{Cash book}{Sales book}{Pass book}
\mcq{118}{The register which is maintained by consignee to keep full record of goods received is called:}{Consignment outward book}{Consignment inward book}{Consignment purchases book}{Consignment in transit book}
\mcq{119}{In Consignment all the expenses, either paid by the consignor himself or by the consignee are borne by the:}{Consignor}{Consignee}{Purchases}{All of these}
\mcq{120}{Consignee:}{Bear express or loss in consignment}{Share profit made on the sale of goods}{Receives his commission from the consignor}{All of the above}
\mcq{121}{In case of Consignment, risk attached to the goods sold before the goods consigned are sold lies with the:}{Consignee}{Consignor}{Buyer}{Debtor}
\mcq{122}{The remuneration of the consignee for selling the goods of the consigner is called:}{Salary}{Interest}{Dividend}{Commission}
\mcq{123}{It is paid to consignee if he takes the risk of bad debts:}{Commission}{General commission}{Ordinary commission}{Delcredre commission}
\mcq{124}{The relationship between seller and buyer of goods is that of a:}{Principal and agent}{Debtor and creditor}{Banker and customer}{Consignor and customer}
\mcq{125}{In consignment expenses after delivery of goods to consignee are borne by:}{Purchaser}{Seller}{Consignor}{Consignee}
\mcq{126}{In sale expenses after delivery of goods are borne by:}{Purchaser}{Seller}{Consignor}{Consignee}
\mcq{127}{In consignment the risk attached to goods delivered is with:}{Seller}{Buyer}{Consignor}{Consignee}
\mcq{128}{In sale the risk and damages attached to goods sold are transferred to:}{Seller}{Buyer}{Consignor}{Consignee}
\mcq{129}{Account sales is submitted by:}{Consignee to consignor}{Debtors to creditor}{Principal to agent}{“a” and”b”}
\mcq{130}{In case ofdel-creditedcommission the liability for bad debts will be on:}{Consignee}{Consignor}{Insurance}{None of these}
\mcq{131}{It is paid to agent to work hard to push a new line of product in the market:}{Commission}{Del-credited commission}{Overriding commission}{Ordinary account}
\mcq{132}{In consignment it is prepared to find out the value of profit or loss:}{Consignee account}{Consignor account}{Consignment account}{Goods sent on consignment account}
\mcq{133}{Consignment account is a combination of:}{Trading account and balance sheet}{Trading account and receipts and payments account}{Trading account and statement of affairs}{Trading account and profit and loss account}
\mcq{134}{Abnormal loss occurs due to:}{Loss by fire}{Accident}{Theft}{All of these}
\mcq{135}{Norma loss occurs due to:}{Leakage}{Loss in weight due to nature}{Accident}{“a” and “b”}
\mcq{136}{Forwarding letter send by consignor to consignee along with the goods consigned:}{Invoice}{Performa invoice}{Bill of exchange}{Voucher}
\mcq{137}{The forwarding letter send by consignor to consignee is a substitute of:}{Voucher}{Journal}{Invoice}{Consignment account}
\mcq{138}{It shows the detail about sale of goods, expenses paid by consignee etc:}{Account sales}{Sales account}{Sales book}{Purchases account}
\mcq{139}{Consignment account is a nature of:}{Personal account}{Real account}{Nominal account}{“a” and “b”}
\mcq{140}{Goods sent on consignment is a nature of:}{Personal account}{Real account}{Nominal account}{“a” and “b”}
\mcq{141}{Consignee account is the nature of:}{Real account}{Property account}{Personal account}{None of these}
\mcq{142}{Consignee account is the nature of:}{Real account}{Property account}{Personal account}{None of these}
\mcq{143}{In book of consignor, on receiving advance in cash from consignee should be debited to:}{Cash account}{Consignment account}{Consignee’s account}{None of these}
\mcq{144}{In book of consignee on payment of advance to consignee should be debited to:}{Consignee’s account}{Consignor’s account}{Consignment account}{Cash account}
\mcq{145}{In the book of consignor, if the consignee has accepted a bill of exchange should be debited:}{Bill receivable account}{Bill payable account}{Consignment account}{Consignee’s account}
\mcq{146}{In consignee’s book, the acceptance given of a bill should be debited to:}{Bill receivable account}{Bill payable account}{Consignee’s account}{Consignor’s account}
\mcq{147}{In consignee’s book, the payment of expenses by consignee should be debited to:}{Consignment account}{Consignor’s account}{Cash account}{Expenses account}
\mcq{148}{In consignee’s book the payment of expenses by consignee personally should be credited to:}{Expenses account}{Consignor’s account}{Cash account}{Consignment account}
\mcq{149}{In book of consignor, the goods sold by consignee for cash should be debited to:}{Consignee’s account}{Debtor’s account}{Cash account}{Consignment account}
\mcq{150}{In book of the goods sold by consignee for credit should be debited to:}{Cash account}{Consignment account}{Consignor’s account}{Sundry debtors account}
\mcq{151}{In book of consignee commission should be credited to:}{Consignor’s account}{Consignee’s account}{Cash account}{Commission account}
\mcq{152}{In book of consignor, commission on consignment should be credited to:}{Commission account}{Cash account}{Consignment account}{Consignee’s account}
\mcq{153}{In book of consignee, Del-credited commission should be debited to:}{Commission account}{Consignor’s account}{Del-credited commission account}{Cash account}
\mcq{154}{In book of consignee, the cash received from sundry debtors should be debited to:}{Consignor’s account}{Cash account}{Sundry debtors account}{Consignee’s account}
\mcq{155}{In book of consignor, the unsold stock with the consignee should be debited to:}{Consignment stock account}{Consignment account}{Stock account}{Consignee’s account}
\mcq{156}{In book of consignor’s abnormal loss of goods should be credited to:}{Loss of goods account}{Profit and loss account}{Consignment account}{Consignment account}
\mcq{157}{In book of consignor, profit on consignment should be credited to:}{Profit and loss account}{Cash account}{Consignee’s account}{Consignment account}
\mcq{158}{The entry for goods sent on consignment in consignee’s book will be:}{Purchase account ..DrTo consignor’s account}{Purchase account…DrTo cash account}{Purchases account …DrTo consignment account}{No entry}
\mcq{159}{The entry for expenses paid by consignor in consignor’s book will be:}{Expenses account …DrTo cash account}{Expenses account …DrTo consignee’s account}{Consignment account …DrTo cash/bank account}{No entry}
\mcq{160}{The entry for expenses paid by consignor in consignee’s book will be:}{No entry}{Expenses account …DrTo cash account}{Expenses account …DrTo consignor’s account}{Expenses account …DrTo consignment account}
\mcq{161}{The entry for expenses paid by consignee in consignment in consignee’s book will be:}{No entry}{Consignor’s account ..DrTo cash/bank account}{Expenses account ..DrTo cash/bank account}{Expenses account ..DrTo consignor’s account}
\mcq{162}{The entry for expenses paid by consignee in consignor’s book will be:}{No entry}{Consignment account …DrTo cash account}{Consignment account …DrTo consignee’s account}{Consignee’s account …Dr To consignment account}
\mcq{163}{The entry for payment of advance in cash by consignee to consignor in consignee’s book will be:}{Consignment account ..Dr To cash account}{Consignment account ..DrTo consignor’s account}{No entry}{Consignor’s account ..DrTo cash account}
\mcq{164}{The entry for acceptance of bill in advance by consignee in consignee’s book will be:}{Bill payable account ..DrTo consignor’s account}{Bill receivable account..DrTo consignor’s account}{No entry}{Consignor’s account ..DrTo bill payable account}
\mcq{165}{The entry for acceptance given of bill in advance by consignee in consignor’s book will be:}{No entry}{Bill receivable account ..DrTo consignee’s account}{Bill receivable account ..DrTo consignment account}{Consignment account ..DrTo consignee’s account}
\mcq{166}{The entry of expense to be borne by consignee personally, the entry in consignee’s book will be:}{No entry}{Expenses account …DrTo cash account}{Expenses account …DrTo consignor’s account}{Consignment account ..DrTo cash account}
\mcq{167}{When goods are sold for cash the entry in consignee’s book will be:}{Cash account … DrTo sales account}{Cash account … DrTo consignor’ s account}{Consignor account …DrTo sales account}{No entry}
\mcq{168}{When the goods are sold on credit by consignee, the entry in consignor’s book will be:}{Sunday debtors account …Dr To sales account}{Cash account ..Dr To sundry debtors account}{Consignee account …DrTo consignment account}{No entry}
\mcq{169}{The entry for commission in consignment in consignee’s book will be:}{Cash account ..Dr To commission account}{Consignor’s account .DrTo commission account}{Consignment account …DrTo consignee’s account}{No entry}
\mcq{170}{In consignment, the entry for commission in consignor’ s book will be:}{Commission account ..DrTo cash account}{Commission account …DrTo consignee’s account}{No entry}{Consignment account ..DrTo consignee’s account}
\mcq{171}{The entry for goods taken over by consignee in consignment personally in consignee’s book will be:}{Goods account …DrTo consignor’s account}{No entry}{Purchase account ..DrTo consignor’s account}{Purchase account ..DrTo cash account}
\mcq{172}{For the goods taken over by consignee for personally use the entry in consignor’s book will be:}{Cash account ..DrTo sales account}{Consignee’s account …DrTo consignment account}{Consignee’s account …DrTo sales account}{Purchase account …DrTo consignee’s account}
\mcq{173}{When cash is received from sundry debtors by consignee, the entry in consignee’s book will be:}{Cash account … DrTo sundry debtors account}{Consignor’s account …DrTo sundry debtors account}{No entry}{Cash account ..DrTo consignor’s account}
\mcq{174}{When cash is received from debtors, by consignee, the entry in consignor’s book will be:}{Cash account … Dr To consignee’s account}{No entry}{Consignee’s account …DrTo consignment account}{Cash account … Dr To sundry debtors account}
\mcq{175}{In consignment, if the consignee does not receive the del credited commission, the entry in consignee’s book for bad debts will be:}{Bad debts account …DrTo sundry debtors account}{No entry}{Consignee’s account …DrTo sundry debtors account}{Bad debts account…DrTo consignor’s account}
\mcq{176}{In consignment, if the consignee receives the del credited commission the entry for bad debts in consignee’s book will be:}{Bad debts account…DrTo sundry debtors account}{No entry}{Bad debts account …DrTo consignor’s account}{Commission account ..DrTo sundry debtors account}
\mcq{177}{In consignment, for unsold stock with consignee’s the entry in consignee’s book will be:}{Consignment stock account … Dr To consignment account}{Consignment stock account … Dr To consignor’s account}{No entry}{Consignment stock account … Dr To trading account}
\mcq{178}{In consignment, in case of abnormal loss of goods, the entry for the value of goods damaged in consignor’s book will be:}{Goods account …DrTo consignment account}{Consignment account …DrTo abnormal loss account}{Consignee’s account…DrTo abnormal loss account}{Loss of stock account…DrTo consignment account}
\mcq{179}{In consignment, in case of abnormal loss of goods, the entry for the value of goods damaged in consignee’s book will be:}{Loss on stock account …DrTo consignor’s account}{Profit \& loss account…DrTo loss on stock account}{No entry}{Consignor’s account…DrTo loss on stock account}
\mcq{180}{The claim against abnormal loss which h is accepted by insurance co, the entry in consignor’s book will be:}{Loss on stock account ..DrTo consignment account}{Bank/insurance account…Dr To loss on stock account}{No entry}{Consignment account..DrTo loss on stock account}
\mcq{181}{The claim against abnormal loss which is accepted by insurance company, the entry in consignee’s book will be:}{No entry}{Loss on stock account..DrTo consignment account}{Bank/insurance co account …Dr To loss on stock account}{Profit \& loss account..DrTo loss on stock account}
\mcq{182}{The entry for goods sold which are damaged due to abnormal loss in consignee’s book will be:}{Cash account… DrTo sales account}{No entry}{Cash account … DrTo consignor’s account}{Consignor’s account…DrTo sales account}
\mcq{183}{The entry for goods sold by consignee which are damaged due to abnormal loss in consignor’s book will be:}{No entry}{Consignee’s account …DrTo consignment account}{Loss of stock account…DrTo consignment account}{Consignee’s account…DrTo loss on stock account}
\mcq{184}{The businessmen can increase the sale of their products:}{By opening their own branches in different parts of the countries and abroad}{Byadopting agents in their own country as well all as abroad}{By a huge head office building}{A and B}
\mcq{185}{The entry for goods sent on consignmentin consignee's book will be:}{Purchase account ....Dr to Consignor' account}{Purchase account... Dr to cash account}{Purchase account ..... Dr to consignment account}{No entry}
\mcq{186}{The entry for payment of expenses to be borne by consignee personally in consignor's book will be:}{No entry}{Bill receivable account .....Dr to consignee's account}{DrBill receivable account .... Dr to consignment account}{Consignment account..... Dr to consignee's account}
\mcq{187}{The entry for Del credere commission in consignment in consignee's book will be:}{Cash account ......Dr to commission account}{Cash account.....Dr to Del credere account}{No entry}{Consignor's account ..... Dr to Del credere commission account}
\mcq{188}{In consignment, the entry for net loss of stock due to fire, accident, theft in consignor’s book will be:}{Profit \& loss account …DrTo loss on stock account}{Loss on stock account…Dr.To consignment account}{Consignee’s account…DrTo consignment account}{No entry}
\mcq{189}{In consignment the entry for net loss of stock due to fire, accident theft in consignee’s book will be:}{Profit \& loss account ..DrTo loss on stock account}{No entry}{Loss of stock account …DrTo consignor’s account}{Consignor’s account …DrTo profit and loss account}
\mcq{190}{For final payment in cash or bank draft to consignor by consignee, the entry in consignee’s book will be:}{Consignor’s account ..DrTo cash/bank account}{Cash/bank account…DrTo consignor’s account}{Cash/bank account …DrTo consignee’s account}{Consignee’s account …Dr To consignment account}
\mcq{191}{The entry for final payment by accepting a bill by consignee in consignor’s book will be:}{Consignor’s account …DrTo bill receivable account}{No entry}{Bill payable account…DrTo consignee’s account}{Bill receivable account …Dr To consignee’s account}
\mcq{192}{The entry for final payment by accepting a bill by consignee in consignee’s book will be:}{Consignor’s account …DrTo bill payable account}{No entry}{Bill payable account …DrTo consignor’s account}{Cash account … DrTo bill payable account}
\mcq{193}{The entry for profit on consignment in consignor’s book will be:}{Profit and loss account …Dr To consignee’s account}{Profits and loss account ..Dr To consignment account}{No entry}{Consignment account ..Dr To profit and loss account}
\mcq{194}{The entry for profit on consignment in consignee’s book will be:}{No entry}{Consignment account …Dr To consignor’s account}{Consignor’s account ..DrTo profit and loss account}{Cash account … DrTo profit and loss account}
\mcq{195}{The entry for loss on consignment in consignee’s book will be:}{No entry}{Profits and loss account …Dr To consignor ‘s account}{Consignor’s account …DrTo profit and loss account}{Cash account … DrTo profit and loss account}
\mcq{196}{Joint stock company is formed under:}{1913Finance society act 1913}{1932Partnership act1932}{1984Companies ordinance 1984}{1865Common carrier act 1865}
\mcq{197}{It is an artificial person recognized by law:}{Joint stock company}{Partnership}{Sole proprietorship}{“b” and “c”}
\mcq{198}{The official signature of the company:}{Partnership deed}{Common seal}{Common act}{19841984}
\mcq{199}{Which business has the perpetual existence:}{Sole proprietorship}{Partnership}{“a” and “b”}{Joint stock company}
\mcq{200}{How many kinds of companies are registered under company ordinance 1984:}{Two}{Three}{Four}{Five}
\mcq{201}{This is the principle form of company:}{Company limited by shares}{Company limited by guarantee}{Unlimited company}{Association not for profit}
\mcq{202}{The company in which members are personally liable for the debts of company:}{Company limited by shares}{Unlimited company}{Company limited by guarantee}{Association not for profit}
\mcq{203}{The liability of the members is limited according to undertaking of members to contribution to the assets of the company in the event of its being wound up:}{Company limited by shares}{Company limited by guarantee}{Unlimited company}{Association not for profit}
\mcq{204}{This association is mainly formed for the promotion of commerce, art, religion, charity etc:}{Company limited by shares}{Unlimited company}{Company limited by guarantee}{Association not for profit}
\mcq{205}{The company which can be formed by at least two members:}{Private company}{Public company}{Sole company}{All of the above}
\mcq{206}{The company which can be formed by at least seven members:}{Private company}{Public company}{Partnership}{Sole proprietorship}
\mcq{207}{Maximum members in private company are:}{10}{50}{100}{150}
\mcq{208}{Maximum members in public company are:}{1010}{100100}{150150}{None of these}
\mcq{209}{In private company directors should be least:}{4}{3}{2}{1}
\mcq{210}{In partnership, in case of banking business, the partners must be exceed:}{5}{7}{10}{15}
\mcq{211}{In partnership other than banking business the members must be exceed:}{20}{15}{10}{5}
\mcq{212}{A company is formed by:}{Shareholders}{Promoters}{Directors}{Underwriters}
\mcq{213}{Shareholders are:}{Owners}{Directors}{Managers}{Employees}
\mcq{214}{A company is managed by:}{Shareholders}{Managers}{Directors}{Employees}
\mcq{215}{How many legal basic documents of a public company are:}{Two}{Three}{Four}{Five}
\mcq{216}{It contains the fundamental conditions upon which alone a company can incorporate:}{Articles of association}{Prospectus}{Memorandum of association}{Partnership deed}
\mcq{217}{It contains the rules and regulations for the internal management of company:}{Articles of association}{Partnership deed}{Memorandum of association}{Prospectus}
\mcq{218}{It is a valuable document issued by the company for raising of capital:}{Memorandum of association}{Articles of association}{Statutory report}{Prospectus}
\mcq{219}{Formation of a public limited company is based on:}{Promotion}{Incorporation}{Commencement of business}{All of these}
\mcq{220}{Which one of the following statement is submitted to registrar in the absence of prospectus for getting the certificate of commencement of business in case of public company:}{Income statement}{Statement in lieu of prospectus}{Statement of profit or loss}{Earning statement}
\mcq{221}{He gives the guarantee for shares and debentures offered to public:}{Director}{Promoter}{Underwriter}{Over writer}
\mcq{222}{The remuneration which is paid by company to underwriter is called:}{Commission}{Del-Credit commission}{Overriding commission}{Under-writing commission}
\mcq{223}{A person who purchase the shares of a company becomes its:}{Shareholder}{Director}{Limited}{Unlimited}
\mcq{224}{the liability of a member in a company is generally:}{Current}{Contingent}{Limited}{Unlimited}
\mcq{225}{The shares of a public limited company are generally:}{Freely transferable}{Not transferable}{Restricted}{All of these}
\mcq{226}{A private company can commerce business after its incorporation:}{After two months}{After two weeks}{Immediately}{After three weeks}
\mcq{227}{One of the units into which the capital of the company is divided is called:}{Debenture}{Share}{Capital}{Prize bond}
\mcq{228}{The sum of the par value of the shares of a company is called:}{Share}{Working capital}{Reserve capital}{Share capital}
\mcq{229}{The amount of share capital with which a company is registered is called:}{Authorized capital}{Issued capital}{Paid up capital}{Called up capital}
\mcq{230}{It is stated in memorandum of association and articles of association of the company:}{Debentures}{Shares}{Paid up capital}{Called up capital}
\mcq{231}{Authorized capital is also called:}{Registered capital}{Nominal capital}{Working capital}{“a” and “b”}
\mcq{232}{The face value of a share is also known as:}{Marker value}{Par value}{Book value}{All of these}
\mcq{233}{Shares offered to general public for contribution are called:}{Authorized shares}{Called up shares}{Issued shares}{Subscribed shares}
\mcq{234}{The total par value of shares taken up by the public is called:}{Authorized capital}{Issued capital}{Subscribed capital}{Paid up capital}
\mcq{235}{It is that portion of subscribed capital which has been called up by the company:}{Called up capital}{Paid up capital}{Subscribed capital}{Issued capital}
\mcq{236}{The portion of called up capital which has been actually paid by shareholders is called:}{Called up capital}{Issued capital}{Paid up capital}{Reserve capital}
\mcq{237}{The value of shares according to books of accounts of the company:}{Book value}{Par value}{Nominal value}{Sale value}
\mcq{238}{This is the value of a share as quoted on the stock exchange:}{Book value}{Par value}{Face value}{Market value}
\mcq{239}{The expenses incurred in the initial stages of incorporation:}{Direct expenses}{Indirect expenses}{Primary expenses}{Operating expenses}
\mcq{240}{The entry for share application money received will be:}{Ordinary share application a/c … Dr To Banka/c}{No entry}{Bank a/c …Dr To Ordinary share application a/c}{Bank a/c …Dr To Ordinary share capital a/c}
\mcq{241}{Company does not receive the application for all shares issued in:}{Under subscription}{Over subscription}{Subscription}{Normal subscription}
\mcq{242}{Company receives more application than shares issued in:}{Under subscription}{Over subscription}{Normal subscription}{Routine subscription}
\mcq{243}{The entry for share application money refunded will be:}{Bank a/c …Dr To Ordinary share application a/c}{Bank a/c …Dr To Ordinary share capital a/c}{Ordinary share application a/c … Dr To Bank a/c}{Ordinary share capital a/c …Dr To Bank a/c}
\mcq{244}{The entry for share application money refunded will be:}{Bank a/c …Dr To Ordinary share application a/c}{Bank a/c …Dr To Ordinary share capital a/c}{Ordinary share application a/c … Dr To Bank a/c}{Ordinary share capital a/c …Dr To Bank a/c}
\mcq{245}{If shares are issued at face value by company then it is called:}{Shares issued at par}{Shares issued at premium}{Shares issued at discount}{Shares issued at book value}
\mcq{246}{Issue of shares at a price more than the face value or normal value is called:}{Issue of shares at face value}{Issue of shares at par value}{Issue of shares at discount}{Issue of shares at premium}
\mcq{247}{Issue of shares at a price less than the face value is called:}{Issue of shares at nominal value}{Issue of shares at par value}{Issue of share at discount}{Issue of shares at face value}
\mcq{248}{Premium on issue of shares is a:}{Business loss}{Revenue profit}{Capital profit}{Revenue loss}
\mcq{249}{Premium on issue of shares may be used to write off:}{Preliminary expenses}{Discount on issue of shares}{Commission paid on issue of shares}{All of these}
\mcq{250}{When shares are issued at a price which is less than its par value it is said that it has been issued at:}{Face value}{Par value}{Discount}{Premium}
\mcq{251}{If a share of Rs.10 is issued at 10\% discount, then issue price of share will be:}{Rs 11}{Rs 9}{Rs 10}{Rs 12}
\mcq{252}{If a share of Rs.10 is issued at 10\% premium, then issue price of share will be:}{Rs 9}{Rs 11}{Rs 12}{Rs 10}
\mcq{253}{If a share of Rs.100 is issued at Rs.90 then it is to be issued at:}{5\% discount}{10\% premium}{5\% discount}{10\% premium}
\mcq{254}{If a share of Rs.100 is issued at Rs.110 then it is to be issued at:}{5\% discount}{5\% premium}{10\% discount}{10\% premium}
\mcq{255}{The percentage of discount or premium is calculated on share:}{Par value}{Discount value}{Premium value}{Market value}
\mcq{256}{According to companies ordinance, the maximum rate of discount should be exceed:}{20\%}{15\%}{10\%}{5\%}
\mcq{257}{If the issue of shares at discount is sanctioned by authority then the shares must be issued within:}{30 days}{40 days}{50 days}{60 days}
\mcq{258}{Discount on issue of shares is a:}{Revenue loss}{Capital profit}{Capital loss}{Revenue profit}
\mcq{259}{Discount on issue of shares is written in:}{Trading account}{Credit side of profit and loss account}{Asset side of balance sheet}{Liability side of balance sheet}
\mcq{260}{Premium on issue of shares is shown on:}{Debit side of trading account}{Debit side of profit and loss account}{Asset side of balance sheet}{Liability side of balance sheet}
\mcq{261}{Ordinary shares are also known as:}{Preferred shares}{Equity shares}{Right shares}{Bonus shares}
\mcq{262}{The excess shares application money should be refunded with in:}{10 days}{20 days}{15 days}{5 days}
\mcq{263}{Shareholders are the:}{Owners}{Debtors}{Creditors}{Promoters}
\mcq{264}{Debenture holders are the:}{Debtors of the company}{Creditors of the company}{Owners of the company}{Directors of the company}
\mcq{265}{The value which is printed on shares is called:}{Face value}{Market value}{Book value}{Purchase value}
\mcq{266}{It is a document evidencing indebtedness of the company:}{Share}{Capital}{Debenture}{Bill of exchange}
\mcq{267}{Debenture includes:}{Debenture stock}{Bonds}{Participation term certificates}{All of these}
\mcq{268}{The debentures which do not carry any charge on the assets of the company are called:}{Naked debentures}{Mortgage debentures}{Secured debentures}{None of these}
\mcq{269}{These are repayable at the end of specified period:}{Redeemable debentures}{Irredeemable debentures}{Non-redeemable debentures}{“a” and “b”}
\mcq{270}{These debentures are repayable after a long period of time or never payable during the existence of the company:}{Redeemable debentures}{Payable debentures}{Irredeemable debenture}{“a” and “b”}
\mcq{271}{These debentures can be transferred by mere delivery:}{Order debentures}{Bearer debentures}{Registered debentures}{Crossed debentures}
\mcq{272}{These debentures cannot be transferred by mere delivery:}{Unregistered debentures}{Registered debentures}{Bearer debentures}{Crossed debentures}
\mcq{273}{The conversion of debentures into debenture stock is made by:}{Promoters}{Shareholders}{Directors}{underwriters}
\mcq{274}{The entry for issue of debentures to vendor’s company at par will be:}{Bank a/c …Dr To debentures a/c}{Debentures a/c …Dr To bank a/c}{Vendor’s company a/c ..DrTo debentures a/c}{Vendor’s company a/c … Dr To bank a/c}
\mcq{275}{What will be entry when debentures are issued for cash at par:}{nullBank a/c … Dr To debentures a/c}{nullBank a/c … Dr To ordinary share capital a/c}{nullOrdinary share capital a/c …Dr To bank a/c}{nullDebentures a/c …DrTo bank a/c}
\mcq{276}{What will be the entry when debentures are issued for each at discount:}{nullBank a/c … DrTo debentures a/c}{nullBank a/c … Dr To discount on issue of debentures a/cTo debentures a/c}{nullBank a/c … Dr To Discount on issue of debentures a/cTo debentures a/c}{nullBank a/c … DrTo discount on issue of debentures a/c}
\mcq{277}{What will be the entry when debentures are issued as collateral security:}{nullNo entry}{nullBank a/c … Dr To debentures suspense a/c}{nullDebentures a/c …Dr To bank loan a/c}{nullDebentures suspense a/c … Dr To debentures a/c}
\mcq{278}{Ifdebentures are issue at discount and redeemable at par than the entry at the time of redemption will be:}{nullBank A/C...... Dr to Loss on issue ofdebentures A/C .....Dr to Premium on redemption ofdebentures}{nullDebentures A/C..... Dr to premium on redemption ofdebentures A/C Bank A/C}{nullBank A/C Discount on issue ofdebentures A/C}{nullBank A/Cdebentures A/C premium on redemption ofdebentures A/C}
\mcq{279}{Ifdebentures are issued at discount at redeemable at premium then the entry at the time of issue will be:}{nullBank A/C .....Dr to loss on issue ofdebentures A/Cdebentures A/C premium on redemption ofdebentures A/C}{nullDebentures A/C ........ Dr to premium on redemption ofdebentures A/C Bank A/C}{nullDebentures A/C......Dr to Bank A/C}{nullBank A/C premium on redemption ofdebentures A/C ......Dr todebentures A/C}
\mcq{280}{Ifdebentures are issued at par and redeemable at premium then the entry at the time of redemption will be:}{nullDebentures A/C ........ Dr to premium on redemption ofdebentures A/C debentures A/C}{nullBank A/C .....Dr to loss on issue ofdebentures A/Cdebentures A/C premium on redemption ofdebentures A/C}{nullDebentures A/C ........ Dr to premium on redemption ofdebentures A/C Bank A/C}{nullBank A/C .....Dr to loss on issue ofdebentures A/Cdebentures A/Cdebentures A/C}
\mcq{281}{If debentures are issued at par and redeemable at premium then the entry at the time of issue will be:}{nullBank A/C .....Dr to loss on issue ofdebentures A/Cdebentures A/C premium on redemption ofdebentures A/C}{nullDebentures A/C ........ Dr to premium or redemption ofdebentures A/C Bank A/C}{nullBank A/C premium on redemption ofdebentures A/C ......Dr todebentures A/C}{nullDebentures A/C ........ Dr to premium or redemption ofdebentures A/C}
\mcq{282}{If debentures are issue at premium and redeemable at par then the entry at the time of redemption will be:}{nullBank A/C ... Dr to premium on issue of debentures A/C}{nullDebentures A/C......Dr to Bank A/C}{nullBank A/C .... dr to premium on issue of debentures A/C Debentures A/C}{nullBank A.C premium on redemption of debentures A/C Debentures A/C}
\mcq{283}{If debentures are issue at discount and redeemable at par then the entry at the time of redemption will be:}{nullBank A/C... dr to discount on the issue of debentures A/C Debentures A/C}{nullDebentures A/C......Dr to Bank A/C}{nullBank A/C discounton redemption ofdebentures A/C ......Dr todebentures A/C}{nullBank A/C.... dr to debentures A/C Discount on redemption of debenture A/C}
\mcq{284}{What will be the journal entry at the time of redemption when debentures are issued at par and redeemable at par:}{nullBank A/C ..... dr to Debentures A/C}{nullDebentures A/C ...dr to bank A/C}{nullBank Ac ...... dr to discount on issue of debentures A/C}{nullBank A/C ..... dr to premium on redemption of debentures A/C}
\mcq{285}{What will be the entry when debentures are issued for cash at premium:}{nullBank A/C ... dr to Share premium A/C}{nullBank A/C ..... dr to ordinary share premium A/C Debentures A/C}{nullBank A/C .... dr to premium on issue of debentures A/C Debentures A/C}{nullPremium on issue of debentures A/C ...... dr to Debentures A/C}
\mcq{286}{The entry for issue of debentures of vender's company at discount will be:}{nullBank A/C .... dr to Discount on issue of debentures A/C Debentures A/C}{nullVender's company A/C .... dr to Discount on issue of debentures A/C Debentures A/C}{nullVender's company A/C ..... dr to discount of debentures A/C Debentures A/C}{nullVender's company A/C ..... dr to Debentures A/C}
\mcq{287}{The entry for purchase of business more then the net assets value will be:}{nullSundry assets A/C ... dr to Vender's company A/C}{nullSundry assets A/C .... dr to Vender's company A/C sundry Liabilities}{nullSundry assets A/C .... dr to Goodwill A/C Vender's company A/C Sundry liabilities A/C}{nullSundry assets A/C ... dr to vender's company A/C Sundry liabilities A/C Capital profit A/C}
\mcq{288}{The entry of purchase of business is less than the net assets value will be:}{nullSundry assets A/C... dr to Vendor's company A/C}{nullSundry assets A/C ..... dr to Vender's company A/C Sundry liabilities A/C Capital profit A/C}{nullSundry assets A/C ..... dr to Goodwill A/C vendor's company A/C}{nullSundry assets A/C...... dr to Vendor's company A/C Cash A/C}
\mcq{289}{Depreciation is an:}{Expense}{Revenue}{Liability}{Capital}
\mcq{290}{The gradual decrease in the efficiency of an asset is called:}{Appreciation}{Depreciation}{Construction}{Reconstruction}
\mcq{291}{The permanent and continuous diminution in the quality, quantity or value of an asset:}{Appreciation}{Diminution}{Reduction}{Depreciation}
\mcq{292}{Depreciation:}{Exhaustion of the effective life of an asset}{Diminution in intrinsic value of an asset}{Reduction in the value of fixed asset due to use in business}{All of the above}
\mcq{293}{Internal causes of depreciation:}{Wear and tear}{Depletion}{Obsolescence}{“a” and “b”}
\mcq{294}{Depreciation which occurs due to certain inherent normal causes is known as:}{Internal depreciation}{External depreciation}{Seasonal depreciation}{All of these}
\mcq{295}{The physical deterioration in asset due to use in business:}{Depletion}{Obsolescence}{Wear and tear}{Accident}
\mcq{296}{Decline in the value proportionate to the quantum of production e.g mine, quarry etc:}{Amortization}{Depletion}{Depreciation}{Wear and tear}
\mcq{297}{External causes of depreciation:}{Obsolescence}{Efflux of time}{Accident}{All of these}
\mcq{298}{Decrease in the value of an asset to change in habit or taste of people:}{Efflux of time}{Accident}{Obsolescence}{Wear and tear}
\mcq{299}{The need for provision of depreciation is necessary to:}{Ascertainment of true profit or loss}{Ascertainment of true cost of production}{Replacement of assets}{All of these}
\mcq{300}{We write depreciation on the:}{Debit side of trading account}{Credit side of trading account}{Debit side of profit and loss account}{Credit side of profit and loss account}
\mcq{301}{If depreciation is not taken into account, the value of asset in the books:}{May be more}{May be less}{No change}{May be more or less}
\mcq{302}{According to provisions of companies ordinance 1984 dividend cannot be declared without:}{Appreciation}{Depreciation}{Replacement of assets}{“a” and “b”}
\mcq{303}{Depreciation:}{Reduces productive activity}{Reduces value of asset gradually}{Is a regular loss}{All of these}
\mcq{304}{Fluctuation:}{Dos not reduce productive capacity}{May rise or fall the value of an asset}{Is generally irregular}{All of these}
\mcq{305}{Depreciation is charged on:}{Fixed assets only}{Current assets only}{Both fixed assets and current assets}{All of these}
\mcq{306}{Depreciation occurs:}{Up to one year}{Up to two years}{Till the last day of the estimated working life of asset}{Up to the replacement of old asset}
\mcq{307}{Depreciation is a charge against:}{Expense}{Revenue}{Liability}{Capital}
\mcq{308}{Depreciation may be ascertained exactly in case of:}{Lease holds property}{Machinery}{Patent right}{“a” and “b”}
\mcq{309}{Depreciation of an asset should be exceeds:}{Original cost}{Depreciation value}{Market priceMarket price}{Scrap value}
\mcq{310}{We should deduct the amount of depreciation from:}{Capital}{Fixed assets}{Cash}{Liability}
\mcq{311}{The decrease in the value of intangible assets is known as:}{Amortization}{Depreciation}{Appreciation}{Depletion}
\mcq{312}{The value of an asset at the end of working life is called:}{Original cost}{Book value}{Scrap value}{Market value}
\mcq{313}{All expenses involved in carrying and installing the asset to the site:}{Cost price of an asset}{Sale price of an asset}{Market price of an asset}{Scrap value of an assetScrap value of an asset}
\mcq{314}{The period during which the asset will help in earning income of business is known as:}{Consumed life}{Expired life}{Exhausted life}{Working life}
\mcq{315}{The decrease in the value of trade mark, copy right, patent right etc is termed as:}{Depreciation}{Amortization}{Depletion}{Reconstruction}
\mcq{316}{The decrease in the value of mines, quarries etc is termed as:}{Depreciation}{Construction}{Depletion}{Amortization}
\mcq{317}{The term depreciation is used with reference to:}{Tangible assets}{Intangible assets}{Current assets}{Fixed assets}
\mcq{318}{The change in shape of an asset due to friction, movement etc. is known as:}{Accident}{Wear and tear}{External cause}{Obsolescence}
\mcq{319}{The decrease in the value of an asset on account of sheer passage of time is:}{Internal depreciation}{External depreciation}{Wear and tear}{Depletion}
\mcq{320}{If the loss on account of depreciation is not considered at the year end, profit will be:}{Less}{More}{No change}{“a” and “b”}
\mcq{321}{Depreciation arises because of:}{Fall in the market value of an asset}{Physical wear and tear}{Fall in value of money}{Increase in the value of money}
\mcq{322}{Due to fluctuation in the market prices, the value of an asset:}{May increase}{May fall}{“a” and “b”}{May not change}
\mcq{323}{Which one of the following asset is to be appreciated in value:}{Land}{Building}{Furniture}{Machinery}
\mcq{324}{The method in which depreciation of an asset is equal in each year:}{Diminishing balance method}{Reducing balance method}{Fixed installment method}{All of these}
\mcq{325}{Scrap value is also known as:}{Resident value}{Break up value}{“a” and “b”}{None of these}
\mcq{326}{In fixed installment method, deprecation is calculated on:}{Book value}{Original cost}{Market price}{“a” and “b”}
\mcq{327}{Original cost of asset Rs45,000; scrap value Rs.3000; useful life of asset 7 years, then the annual value of depreciation will be:}{Rs.6000}{Rs.7000}{Rs.6500}{Rs.8000}
\mcq{328}{Under straight line method, the amount of depreciation:}{Decreases every year}{Constant every year}{Increases every year}{“a” and “b”}
\mcq{329}{The journal entry for depreciation will be:}{nullDepreciation account …Dr To asset account}{nullDepreciation account …DrTo cash account}{nullDepreciation account …DrTo bank account}{nullAsset account …DrTo depreciation account}
\mcq{330}{The entry for transferring the amount of depreciation to profit loss account will be:}{nullDepreciation account …DrTo trading account}{nullDepreciation account …DrTo profit or loss account}{nullProfit and loss account … Dr To depreciation account}{nullProfit and loss account … Dr To asset account}
\mcq{331}{The entry for sale of scrap of an asset will be:}{nullCash account ….Dr To asset account}{nullCash account …Dr To scrap account}{nullAsset account …Dr To cash account}{nullScrap account …Dr To cash account}
\mcq{332}{What will be entry when profit is earned on sale of asset:}{nullCash account …Dr To profit and loss account}{nullProfit and loss account … Dr To cash account}{nullAsset account …Dr To profit and loss account}{nullProfit and loss account … Dr To asset account}
\mcq{333}{What will be the entry for loss on sale of asset:}{nullCash account …Dr To profit and loss account}{nullProfit and loss account … Dr To cash account}{nullAsset account …Dr To profit and loss account}{nullProfit and loss account … Dr To asset account}
\mcq{334}{If the original cost of asset Rs.10,5000; scrap value Rs. 5000; useful life 5 years, then the percentage of depreciation will be:}{10\%}{15\%}{20\%}{25\%}
\mcq{335}{The fixed installment method is also known as:}{Reducing balance method}{Straight line method}{Original cost method}{“b” and “c”}
\mcq{336}{Under reducing balance method the depreciation is always changed on:}{Original cost}{Book value}{Breakup value}{Residual value}
\mcq{337}{Reducing balance method is also known as:}{Diminishing balance method}{Book value method}{Written down value method}{All of these}
\mcq{338}{Suppose cost of asset is Rs.1000 and rate of depreciation 10\% p.a then the book value of asset after two years will be:}{Rs.1000}{Rs.810}{Rs.900}{Rs.729}
\mcq{339}{Under diminishing balance method amount of annual depreciation gradually:}{Decrease}{Increase}{Constant}{No change}
\mcq{340}{Diminishing balance method is suitable for those assets which have:}{Long life}{Short life}{Average life}{Temporary life}
\mcq{341}{Which one of the demerit of reducing balance method:}{It takes a very long time to write off an asset to approximately nil}{Rate of depreciation is high}{“a” and “b”}{None of the above}
\mcq{342}{This method is used by income tax authorities for granting depreciation allowance to taxpayers:}{Straight line method}{Written down value method}{Fixed installment method}{Original cost method}
\mcq{343}{Suppose residual value is Rs.64000; original cost Rs.1000000; estimated useful life is 3 years, then the rate of depreciation under diminishing balance method will be:}{80\%}{60\%}{40\%}{20\%}
\mcq{344}{Under reducing balance method depreciation is calculated on:}{Debit balance}{Credit balance}{Cost+depreciation}{“a” and “b”}
\mcq{345}{The amount of depreciation charged on a machinery will be debited to:}{Machinery account}{Cash account}{Depreciation account}{Trading account}
\mcq{346}{If the original cost of an asset is Rs.10000; rate of depreciation 10\% then the value of depreciation under diminishing balance method after second year will be:}{Rs.900}{Rs.1000}{Rs.1100}{Rs.1200}
\mcq{347}{If the original cost of an asset is Rs.10000; rate of depreciation 10\% then the value of depreciation under diminishing balance method after second year will be:}{Rs.1000}{Rs.900}{Rs.800}{Rs.810}
\mcq{348}{If original cost of an asset is Rs.10000, rate of depreciation 10\%, then the value of depreciation under diminishing balance method after third year will be:}{Rs.1000}{Rs.900}{Rs.700}{Rs.810}
\mcq{349}{If original cost of an asset is Rs.10000, rate of depreciation 10\%, then value of depreciation under straight line method after third year will be:}{Rs.1000}{Rs.900}{Rs.800}{Rs.810}
\mcq{350}{The book value of machinery on January 1,2014 is Rs.20000. two years later the book value is Rs.10000. the straight line depreciation rate of charge each year is:}{5\%}{17\%}{25\%}{33\%}
\mcq{351}{Which one is the merit of fixed installment method:}{Simplicity}{Asset can be written of to zero}{Useful for providing depreciation on leasehold property, patents and other like assets}{All of above}
\mcq{352}{Which of the demerit of fixed installment method:}{No provision is made for interest on amount invested in the purchase of asset}{Depreciation rate is high}{Depreciation of asset decrease but the amount of depreciation remains the same}{A and C}
\mcq{353}{Which one of the wasting asset:}{Land}{Goodwill}{Mines}{Patent}
\mcq{354}{Which one of the tangible asset:}{Goodwill}{Trade mark}{Copy right}{Machinery}
\mcq{355}{Which one of the intangible asset:}{Goodwill}{Trade mark}{Copy right}{All of these}
\mcq{356}{The relation between persons who have agreed to share the profits of a business carried on by all or any one of them acting for all is called:}{Sole proprietorship}{Partnership}{Joint stock company}{Individual}
\mcq{357}{The partnership act is:}{19321932}{19121912}{19621962}{19841984}
\mcq{358}{The essential element of partnership is:}{There must be an agreement entered into by all the persons concerned}{The agreement must be to share the profits of a business}{The business must be carried on by all or any of them acting for all}{All of the above}
\mcq{359}{It sets out the terms on which partners agreed to form a partnership:}{Partnership agreement}{Partnership deed}{“a” and “b”}{None of these}
\mcq{360}{Partnership agreement contains:}{Name of partnership business}{Nature of partnership business}{Ratio in which profits and losses are to be shared by partners}{All of the above}
\mcq{361}{A partner has a right to:}{Take part in the conduct of business}{Inspect the books of the firm}{Share in profit}{All of these}
\mcq{362}{In the absence of an agreement, profit and loss are divided by partners in the ratio of:}{Capital}{Time devoted by each partner}{Equally}{Interest on capital}
\mcq{363}{In the absence of an agreement, interest on loan advanced by a partner to the firm is allowed at the rate of:}{Six per cent}{Five per cent}{Twelve per cent}{Seven per cent}
\mcq{364}{In the absence of an agreement, partner shall:}{Be paid salaries}{Not be paid salaries}{Be paid salaries to those who work for the firm}{Paid salaries to all the partners}
\mcq{365}{Current accounts of the partners should be opened when the capital are:}{Fluctuating}{Fixed}{Either fixed or fluctuating}{Variable}
\mcq{366}{Where a partner is entitled to interest on capital subscripted by him, such interest will be payable:}{Only out of profits}{Only out of capitals}{Out of profits or out of capitals}{None of these}
\mcq{367}{The interest on partners’ capital accounts is to be credited to:}{Partners’ capital accounts}{Profit and loss account}{Interest account}{Trading account}
\mcq{368}{In the absence of an agreement to the contrary, the partners:}{Are entitled for 6\% interest on their capital, only when there are profits}{Are entitled for 9\% interest on their capital, only when there are profits}{Are entitled for interest on capital at the bank rate, only when there are profits}{Are not entitled for any interest on their capitals}
\mcq{369}{When interest is allowed on the capitals of the partners, it is calculated on the:}{Capital in the beginning of the year}{Capital at the end of the year}{Capital at the end less drawings if any}{Average capital}
\mcq{370}{The current account of a partner:}{Will always have a credit balance}{Will always have a debit balance}{May have a debit balance or a credit balance}{None of these}
\mcq{371}{In the absence of an agreement, interest on drawings is to be charged by the firm at the rate of:}{5\% p.a5\% p.a}{6\% p.a6\% p.a}{8\% p.a8\% p.a}{None of these}
\mcq{372}{In the absence of an agreement, the partners are entitled to interest on capital at the rate of:}{5\% p.a5\% p.a}{6\%p.a6\%p.a}{10\% p.a10\% p.a}{None of these}
\mcq{373}{In the absence of an agreement, interest on drawings is to be charged by the firm at the rate of:}{5\% p.a5\% p.a}{2\% p.a2\% p.a}{4\% p.a4\% p.a}{None of these}
\mcq{374}{The maximum amount which each partner will draw from the business in anticipation of the share of profit is his:}{Drawings}{Profit}{Loss}{Capital}
\mcq{375}{The amount of loan advanced by a partner to partnership is transferred to:}{Partner’s capital account}{Partner’s drawings account}{Partner’s loan account}{None of these}
\mcq{376}{The entry for interest on drawings will be:}{nullInterest on drawing .....Dr}{nullPartner's capital account..... Dr to cash account}{nullPartner's capital account ...... Dr to interest on drawing account}{nullInterest on drawing account}
\mcq{377}{The debit balance of interest on loan account is transferred to:}{Partner’s capital account}{Profit \& loss account}{Partner’s loan account}{Trading account}
\mcq{378}{When capital accounts are fixed, then all adjustments are made in:}{Partner’s capital account}{Partner’s current account}{Partner’s fixed account}{Partner’s nominal account}
\mcq{379}{When capital accounts are fluctuating, then all adjustments are made in:}{Partner’s capital account}{Partner’s loan account}{Partner’s fixed account}{Partner’s current account}
\mcq{380}{Partner’s current accounts should be opened when capital accounts are:}{Fixed}{Fluctuating}{Current}{Nominal}
\mcq{381}{When capital is introduced by a partner, we should be debited to:}{Cash account}{Capital account}{Current account}{Fixed account}
\mcq{382}{The amount of capital introduced by a partner, we should be credited to:}{Cash account}{Capital account}{Current account}{Fixed account}
\mcq{383}{The interest on capital account in partnership is to debited to:}{nullInterest on capital a/c}{nullInterest on drawings a/c}{nullPartner’s capital account}{nullPartner’s loan a/c}
\mcq{384}{The interest on capital account in partnership is to credited to:}{nullInterest on capital a/c}{nullProfit \& loss a/c}{nullPartner’s capital account}{nullCash a/c}
\mcq{385}{What will be the closing entry for interest on capital account:}{nullInterest on capital account…..Dr To Cash account}{nullInterest on capital account …… Dr To Partner’s capital account}{nullInterest on capital account ……. Dr To Profit \& loss account}{nullProfit \& loss appropriation account …..Dr To Interest on capital account}
\mcq{386}{The entry for salary payable to partner will be:}{nullSalary account ... Dr To Cash account}{nullCash account …. Dr To Salary account}{nullSalary account ….. Dr To Partner’s capital account}{nullPartner’s capital account ….. Dr To Salary account}
\mcq{387}{The costing entry for salary payable to partner will be:}{nullSalary account …. Dr To Cash account}{nullCash account …. Dr To Salary account}{nullSalary account …. Dr To Profit \& loss app. Account}{nullProfit \& loss to app.account … Dr To Salary account}
\mcq{388}{The entry for drawings in cash by partner will be:}{nullDrawing account … Dr To Cash account}{nullCash account …. Dr To Drawing account}{nullPartner’s capital account …..Dr To Cash account}{nullCash account …. Dr To Partner’s capital account}
\mcq{389}{The closing entry for interest on drawings will be:}{nullProfit \& loss appropriation a/c …Dr To Interest on drawings account}{nullProfit \& loss appropriation a/c ….Dr To Partner’s capital account}{nullPartner’s capital account ……Dr To Interest on drawings account}{nullInterest on drawings account …. Dr To Profit \& loss appropriation a/c}
\mcq{390}{The entry for distribution of profit among the partners will be:}{nullProfit \& loss appropriation a/c …. Dr To Cash account}{nullCash account … Dr To Profit \& loss appropriation a/c}{nullProfit \& loss appropriation a/c …. Dr To Partner’s capital account}{nullPartner’s capital account ….Dr To Profit \& loss appropriation a/c}
\mcq{391}{The entry for distribution of loss among the partners will be:}{nullProfit \& loss appropriation a/c …Dr To Partner’s capital account}{nullPartner’s capital account …. Dr To Profit \& loss appropriation a/c}{nullPartner’s capital account … Dr To Cash account}{nullCash account ….Dr To Partner’s capital account}
\mcq{392}{Suppose the profit of partnership before charging manager’s commission is Rs.10,560, then the amount of manager’s commission will @5\% on profit:}{Rs.528}{Rs.428}{Rs.628}{Rs.560}
\mcq{393}{Suppose the profit of partnership after charging manager’s commission is Rs.10560, then the amount of manager’s commission will be @5\% on profit:}{Rs.503}{Rs.506}{Rs.500}{Rs.603}
\mcq{394}{Mr. A drew cash Rs 1000 from business for 3 months, the amount of interest on drawings @6\% p.a will be:}{Rs.25}{Rs.5}{Rs.15}{Rs.10}
\mcq{395}{The benefit and advantage of the good name or reputation of a business:}{Liability}{Goodwill}{Amortization}{Depreciation}
\mcq{396}{The valuation of goodwill arises on:}{The admission of a new partner}{The retirement or death of a partner}{The dissolution or sale of firm’s business}{All of the above}
\mcq{397}{The value of goodwill can be calculated with the help of:}{Average profits method}{Super profits method}{Capitalization method}{All of the above}
\mcq{398}{The value of goodwill can be calculated with the help of:}{Average profits method}{Super profits method}{Capitalization method}{All of the above}
\mcq{399}{The value of goodwill is calculated by multiplying the average of a given number of past years by an agreed number under:}{Average profits method}{Super profits method}{Capitalization method}{Revaluation method}
\mcq{400}{What will be the value of goodwill on two years purchase of the average profit of the last three years The profit earned during the last three years are Rs.10,000, Rs.12,000 and Rs.1,4000:}{Rs.24,000}{Rs.20,000}{Rs.12,000}{Rs.36,000}
\mcq{401}{The excess of actual profits over normal profits is known as:}{Profit}{Loss}{Super profit}{Super loss}
\mcq{402}{Suppose the average profit of last five years is Rs.7,000, normal profits Rs.4,000, then the value of goodwill on the basis of three years purchase under super profits method will be:}{Rs.6,000}{Rs.9,000}{Rs.12,000}{Rs.15,000}
\mcq{403}{Suppose the actual net profit of a business is Rs.7,500, normal rate of return 10\%, the net tangible assets Rs 50,000, then the value of goodwill will be:}{Rs 5,000}{Rs.15,000}{Rs.25,000}{Rs.35,000}
\mcq{404}{Goodwill is valued when:}{Profit is earned}{Loss is suffered}{One year is completed}{New partner is admitted}
\mcq{405}{A new partner may be admitted to a partnership:}{With the consent of all the old partners}{With the consent of any one of the partners}{With the consent of two third of the old partners}{Without consent of old partners}
\mcq{406}{X,Y and Z are partners in a firm, if B is to be admitted as a new partner:}{Old partnership has to be dissolved}{Old firm has to be dissolved}{Both old firm and partnership have to be dissolved}{Neither firm nor partnership needs to be dissolved}
\mcq{407}{The share of profit to a new partner in the future profits without any express agreement as to who will contribute to new partner’s share of profit is implied that old partners contribute:}{Equally}{In proportion of their capital}{In their profit sharing ratio}{None of these}
\mcq{408}{Old profit sharing ratio minus new profit sharing ratio is equal to:}{Sacrificing ratio}{Gaining ratio}{Distributing ratio}{None of these}
\mcq{409}{When the incoming partner pays the firm for goodwill in cash, the amount should be debited in firms books to:}{Goodwill account}{Capital account of the incoming partner}{Cash account}{All of the above}
\mcq{410}{On the admission of a partner if the goodwill account is to be raised this should be debited to:}{Partners’ capital accounts}{Goodwill account}{Cash account}{New partner capital account}
\mcq{411}{At the time of admission of a new partner, goodwill raised should be written off in:}{New profit sharing ratio}{Old partner sharing ratio}{Sacrificing ratio}{Gaining ratio}
\mcq{412}{If the incoming partner is to bring his share of goodwill in cash and there exists any balance in goodwill account, then this goodwill account is to be written off among old partner in:}{New profit sharing ratio}{Old profit sharing ratio}{Sacrificing ratio}{Gaining ratio}
\mcq{413}{On the admission of a new partner the decrease in the value of assets is debited to:}{Profit and loss adjustment account}{Assets account}{Old partners’ capital account}{New partner capital account}
\mcq{414}{On the admission of a new partner the increase in the value of assets is debited to:}{Profit and loss adjustment account}{Assets account}{Old partners’ capital account}{New partner capital account}
\mcq{415}{If the adjustment in the value of assets under taken at the time of admission of a partner shows a profit, it should be credited to the capital accounts of:}{The old partners in the new profit sharing ratio}{All partners in their new profit sharing ratio}{The old partners in their old profit sharing ratio}{The old partners in their gaining ratio}
\mcq{416}{At the time of admission of a partner, undistributed profit appearing in the balance sheet of the old firm should be transferred to capital account of:}{Old partners in the old profits sharing ratio}{Old partners in the new profit sharing ratio}{All the partners in the new profit sharing ratio}{None of the above}
\mcq{417}{Goodwill is an:}{Tangible asset}{Intangible asset}{Wasting asset}{Fictitious asset}
\mcq{418}{Goodwill is an:}{Expense}{Profit}{Asset}{Liability}
\mcq{419}{A business running into losses will have generally no:}{Machinery}{Sundry debtors}{Sundry creditors}{Goodwill}
\mcq{420}{The extra amount charged from a new partner over and above the capital is for:}{Purchase of machinery}{Goodwill}{Purchase of furniture}{Payment of liabilities}
\mcq{421}{On the admission of a new partner in partnership:}{Old partnership has to be dissolved}{Old firm has to be dissolved}{“a” and “b”}{None of the above}
\mcq{422}{In the absence of an agreement, the share of new partner in partnership will be:}{In the proportion of capital}{Equally}{According to work}{None of the above}
\mcq{423}{Sacrifice ratio is equal to:}{Old ratio+new ratio}{Old ratio-new ratio}{New ratio-old ratio}{New ratio+old ratio}
\mcq{424}{Goodwill of a firm of A and B is valued at Rs.30,000. It is appearing in the books at Rs.12,000. C is admitted to 1/4th share, what amount is supposed to bring for goodwill:}{3000RS.3,000}{4500Rs.4,500}{7500Rs.7,500}{10،500Rs.10,500}
\mcq{425}{A and B are partner in a business sharing profits in the ratio of 5;3. They admit C as a partner with 1/4 share in the profits which he acquires 3/4 from A and 1/4 from B. He pays Rs.8,000 for his share of goodwill, A and B will be credited by:}{Rs.5,000 and Rs3,000 respectively}{Rs 6,000 and Rs.2,000 respectively}{Rs.4,000 each}{Rs.5,000 each}
\mcq{426}{A and B are partner in a business sharing profits in the ratio of 5;3. They admit C as a partner with 1/4 share in the profits which he acquires 3/4 from A and 1/4 from B. He pays Rs.8,000 for his share of goodwill, A and B will be credited by:}{Rs.5,000 and Rs3,000 respectively}{Rs 6,000 and Rs.2,000 respectively}{Rs.4,000 each}{Rs.5,000 each}
\mcq{427}{The entry for goodwill paid privately by a new partner will be:}{nullCash account ..Dr To goodwill account}{nullNo entry}{nullCash account ..Dr To new partner’s capital account}{nullGoodwill account …DrTo cash account}
\mcq{428}{When goodwill brought in cash by a new partner, the entry will be:}{nullCash account …Dr To goodwill account}{nullGoodwill account …Dr To cash account}{nullGoodwill account …Dr To old partners’ capital accounts}{nullNo entry}
\mcq{429}{When goodwill brought in cash but retained by a new partner, the entry will be:}{nullCash account …Dr To goodwill account}{nullGoodwill account …Dr To new partner’s capital account}{nullNew partner’s capital account …Dr To goodwill account}{nullGoodwill account …DrTo old partner’s capital account}
\mcq{430}{The amount of goodwill brought in cash by a new partner will be credited to old partner in:}{Gaining ratio}{New ratio}{Sacrifice ratio}{Old ratio}
\mcq{431}{The amount of goodwill brought in cash but withdrawn by old partners will be:}{nullCash account ..Dr To goodwill account}{nullGoodwill account …Dr To old partner’s capital accounts}{nullOld partner’s capital account ..Dr To cash account}{nullOld partner’s capital account ….Dr To goodwill account}
\mcq{432}{The entry for goodwill raised but written off will be:}{nullOld partner’s capital account …Dr To goodwill account}{nullNew partner’s capital account …Dr To goodwill account}{nullAll partner’s capital account… Dr To goodwill account}{nullAll partner’s capital account ..Dr To cash account}
\mcq{433}{The entry for goodwill raised but written off will be:}{nullOld partner’s capital account …Dr To goodwill account}{nullNew partner’s capital account …Dr To goodwill account}{nullAll partner’s capital account… Dr To goodwill account}{nullAll partner’s capital account ..Dr To cash account}
\mcq{434}{On the admission of a new partner, if goodwill will brought in cash, it should be debited to:}{Cash account}{Goodwill account}{New partner’s capital account}{Old partners’ capital accounts}
\mcq{435}{Any amount to be written after the admission of a partner is transferred to the capital accounts of all partners in:}{Capital ratio}{New profit sharing ratio}{Old profit sharing ratio}{Sacrifice ratio}
\mcq{436}{Goodwill raised by the partner at the time as admission of a partner will be written off in:}{Gaining}{New profit sharing ratio}{Old profit sharing}{Sacrifice ratio}
\mcq{437}{On the admission of a partner if the goodwill account is to be raised this should be debited to:}{New partners’ capital account}{Goodwill account}{Old partners’ capital accounts}{Cash account}
\mcq{438}{At the time of admission of a new partner in partnership, goodwill raised should be written off in:}{Gaining ratio}{New profit sharing ratio}{Old profit sharing ratio}{Sacrifice ratio}
\mcq{439}{It is prepared to calculate the profit or loss of old partners on the admission of a new partner:}{Revaluation account}{Realization account}{Profit and loss account}{All of these}
\mcq{440}{Revaluation account is also known as:}{Revaluation account}{Profit \& loss adjustment account}{Profit and loss account}{Capital account}
\mcq{441}{Revaluation profit or loss is credited or debited to old partner’s capital accounts in:}{Gaining ratio}{New profit sharing ratio}{Old profit sharing ratio}{Sacrifice ratio}
\mcq{442}{On admission of a partner, the decrease in the value of assets should be debited to:}{Revaluation account}{Realization account}{Assets account}{New partner’s capital account}
\mcq{443}{On admission of a partner, the increase in the value of assets should be debited to:}{Revaluation account}{Realization account}{Assets account}{Partners’ capital accounts}
\mcq{444}{The increase in the value of liabilities on the admission of a new partner should be credited to:}{Revaluation account}{Liabilities account}{Profit and loss account}{Partners’ capital accounts}
\mcq{445}{The decrease in the value of liabilities on the admission of a new partner should be credited to:}{Profit and loss adjustment account}{Liabilities account}{Profit and loss account}{Partners’ capital accounts}
\mcq{446}{Profit on revaluation should be credited to:}{Revaluation account}{Liabilities account}{Assets account}{Old partners’ capital accounts}
\mcq{447}{Profit on revaluation should be credited to:}{Revaluation account}{Liabilities account}{Assets account}{Old partners’ capital accounts}
\mcq{448}{Revaluation loss should be debited to:}{Revaluation account}{All partners’ capital accounts}{New partner’s capital account}{Old partners’ capital accounts}
\mcq{449}{The decrease in the value of building, stock on the admission of a partner should be debited to:}{Building account}{Stock account}{“a” and “b”}{Revaluation a/c}
\mcq{450}{Amount of general reserve at the time of admission of a partner is debited to:}{General reserve account}{Revaluation account}{New partner’s capital account}{Old partner’s capital accounts}
\mcq{451}{Accumulated profit on the admission of a new partner should be credited to:}{Accumulated profit account}{All partner’s capital account}{New partner’s capital account}{Old partner’s capital account}
\mcq{452}{The journal entry for the decrease in the value of assets on admission of a new partner will be:}{nullRevaluation account …Dr To assets account}{nullAssets account ..Dr To revaluation account}{nullCash account ..Dr To asset account}{nullAssets account…Dr To new partner’s capital account}
\mcq{453}{The journal entry for the increase in the value of assets on admission of a new partner will be:}{nullRevaluation account …DrTo assets account}{nullAssets account …Dr To revaluation account}{nullRealization account…DrTo assets account}{nullAssets account …DrTo realization account}
\mcq{454}{The journal entry for the decrease in the value of liability on admission of a new partner will be:}{nullRevaluation account …DrTo liabilities account}{nullProfit \& loss account …DrTo liabilities account}{nullLiabilities account …DrTo revaluation account}{nullLiabilities account ..DrTo profit \& loss account}
\mcq{455}{The journal entry for the increase in the value of liability on admission of a new partner will be:}{nullRevaluation account ..DrTo liabilities account}{nullRevaluation account …DrTo profit \& loss adjustment account}{nullLiabilities account ..DrTo revaluation account}{nullLiabilities account …DrTo profit \& loss account}
\mcq{456}{The journal entry for revaluation profit will be:}{nullRevaluation account …DrTo old partner’s capital accounts}{nullRevaluation account ..DrTo new partner’s capital accounts}{nullOld partner’s capital accounts …Dr To revaluation account}{nullNew partner’s capital account …Dr To revaluation account}
\mcq{457}{Retirement of a partner dissolves:}{Partnership}{Firm}{Business}{“b \& c”}
\mcq{458}{On retirement of a partner in partnership to calculate:}{New profit sharing ratio}{Goodwill}{Revolution of assets and liabilities}{All of these}
\mcq{459}{On retirement of a partner, the share of retiring partner is taken by:}{New partner}{Remaining partners}{Government}{Legal representatives of retiring partner}
\mcq{460}{In case of retirement funds for repayment may be provided out of:}{Firm’s own resources}{Borrowing by the firm}{Additional funds may be bought in by continuing partners}{All of the above}
\mcq{461}{Suppose A, B \& C are partners, sharing profits in the ratio of 1/2, 1/3, 1/6 respectively. What will be the new ratio if A retires in the absence of an agreement:}{1/3 ; 1/6}{1/3 ; 1/2}{1/3 ;1/5}{1/2 ; 1/6}
\mcq{462}{Suppose A, B \& C are partners, sharing profits in the ratio of 1/2, 1/3, 1/6 respectively. What will be the new ratio if A retires in the absence of an agreement:}{1/3 ; 1/6}{1/3 ; 1/2}{1/3 ;1/5}{1/2 ; 1/6}
\mcq{463}{Gaining ratio is equal to:}{New ratio + Old ratio}{New ratio – Old ratio}{Old ratio – New ratio}{Old ratio + New ratio}
\mcq{464}{Gaining ratio is calculated on:}{Admission}{Retirement}{Death}{‘b’ \& ‘c’}
\mcq{465}{In case of retirement, when goodwill is raised at full value, is credited to:}{Retiring partners’ capital A/c}{Remaining partners’ capital A/c}{All partners’ capital A/c}{Goodwill A/c}
\mcq{466}{In case of retirement, when goodwill is raised but written off is credited to:}{Retiring partners’ capital A/c}{Remaining partners’ capital A/c}{All partners’ capital A/c}{Goodwill A/c}
\mcq{467}{In case of retirement, when goodwill is raised with retiring partner’s share is credited to:}{Retiring partners’ capital A/c}{Cash A/c}{Goodwill A/c}{None of these}
\mcq{468}{In case of retirement, when goodwill is raised with retiring partner’s share but written off is credited to:}{Retiring partners’ capital A/c}{Cash A/c}{Goodwill A/c}{Remaining patterns capital A/c}
\mcq{469}{In case of retirement, when goodwill is raised with retiring partners share but written off to remaining partners in the:}{Gaining ratio}{Sacrifice ratio}{New ratio}{Capital ratio}
\mcq{470}{In case of retirement of a partner, full goodwill is credited to the accounts of:}{Retiring partner}{Remaining partners}{All the partner}{None of these}
\mcq{471}{The profit on the revaluation of assets and liabilities on the retirement of a partner of a should be credited to the capital account / accounts of:}{Retiring partner}{Remaining partners}{All the partners}{None of these}
\mcq{472}{In case of retirement, revaluation loss should be debited to:}{Retiring partners’ capital A/c}{Remaining partners’ capital A/c}{All partners’ capital A/c}{Revaluation A/c}
\mcq{473}{In case of retirement amount of general reserve or any other profit is credited to all the partners in:}{Old profit sharing ratio}{New profit sharing ratio}{Sacrifice ratio}{Gaining ratio}
\mcq{474}{In case of retirement, the entry for goodwill at full value in the books should be:}{nullGoodwill A/c….Dr To Retiring partner’s capital A/c}{nullGoodwill A/c….Dr To Remaining partner’s capital A/c}{nullGoodwill A/c….Dr To partner’s capital A/c}{nullAll Partner’s capital A/c…..Dr To Goodwill A/c}
\mcq{475}{The entry for goodwill raised but written off, in case of retirement should be:}{nullRetiring partner’s capital A/c….Dr To Goodwill A/c}{nullRemaining partner’s capital A/c….Dr To Goodwill A/c}{nullGoodwill A/c….DrTo All Partner’s capital A/c}{nullAll Partner’s capital A/c….Dr To Goodwill A/c}
\mcq{476}{The entry for goodwill raised with retiring partner’s share, in case of retirement should be:}{nullGoodwill A/c….Dr To Retiring partner’s capital A/c}{nullGoodwill A/c….Dr To Remaining partner’s capital A/c}{nullGoodwill A/c….Dr To All Partner’s capital A/c}{nullRetiring partner’s capital A/c….Dr To Goodwill A/c}
\mcq{477}{The entry for goodwill raised with retiring partner’s share but written off, in case of retirement should be:}{nullRetiring partner’s capital A/c….Dr To Cash A/c}{nullRemaining partner’s capital A/c….Dr To Goodwill A/c}{nullGoodwill A/c….Dr To Retiring partner’s capital A/c}{nullAll Partner’s capital A/c….Dr To Goodwill A/c}
\mcq{478}{The balance of memorandum revaluation account is transferred to capital accounts of all partners in:}{nullCapital ratio}{nullOld profit sharing ratio}{nullNew profit sharing ratio}{nullSacrifice ratio}
\mcq{479}{All accumulated losses are transferred to the capital accounts of the all partners in case of retirement in:}{Capital ratio}{Old profit sharing ratio}{New profit sharing ratio}{Sacrifice ratio}
\mcq{480}{The credit balance of retiring partner’s capital account if not paid in cash should be transferred to:}{Retiring partner’s capital A/c}{Retiring partner’s loan A/c}{Old partners’ capital A/c}{Realization A/c}
\mcq{481}{Interest on retiring partner loan in the absence of an agreement is paid @ rate of:}{3\% p.a3\% p.a}{5\% p.a5\% p.a}{6\% p.a6\% p.a}{8\% p.a8\% p.a}
\mcq{482}{On the death of a partner, the amount of joint life policy should be credited to the capital accounts of:}{All partners}{Remaining partners}{Deceased partners}{Executor}
\mcq{483}{In case of death of a partner, the amount received from insurance company equal to:}{50\%50\% of policy amount}{25\%25\% of policy amount}{75\%75\% of policy amount}{100\%100\% of policy amount}
\mcq{484}{On the death of a partner, the amount received from joint policy should be credited to:}{Deceased partner’s capital accounts in the old profit sharing ratio}{Remaining partners’ capital accounts in the old profit sharing ratio}{All partners’ capital account equally}{All partners’ capital account in the old profit sharing ratio}
\mcq{485}{Joint life policy account after the maturity of the policy should be transferred to the capital accounts of the partners in:}{Capital ratio}{Old profit sharing ratio}{New profit sharing ratio}{Gaining ratio}
\mcq{486}{The amount payable to retiring partner is shown in the balance sheet of partnership as a:}{Capital}{Loan}{Asset}{Investment}
\mcq{487}{Revaluation account is prepared to determine the profit or loss on the revaluation of assets and liabilities when a:}{Partner dies}{Partner retires}{New partner enter}{All of these}
\mcq{488}{At the time of payment to the retiring partner, the following journal entry will be passed:}{nullRetiring partner's capital A/C.....Dr to Cash/Bank/ A/c}{nullCash/Bank A/c.... Dr to Retiring partner's capital A/c}{nullRemaining partners capital A/c ..... Dr to Cas/Bank A/c}{nullNone of these}
\mcq{489}{The end or termination of partnership is called:}{Admission}{Retirement}{Death}{Dissolution}
\mcq{490}{Dissolution of firm by the consent of all the partners is called:}{Dissolution by agreement}{Dissolution by court}{Dissolution on the happening of certain contingencies}{Dissolution by notice}
\mcq{491}{Dissolution due to completion of venture, death of partner, insolvency of partner or expiry of the terms fixed is called:}{Dissolution by agreement}{Dissolution of notice}{Dissolution on the happening of certain contingencies}{Dissolution by court}
\mcq{492}{Dissolution of firm due to the insolvency or due to the happening of an event, which makes it unlawful for partners to carry on partnership is called:}{Dissolution by agreement}{Compulsory dissolution}{Dissolution on the happening of certain contingencies}{Dissolution by court}
\mcq{493}{The court may dissolution the firm:}{When a partner has become of unsound mind}{A partner is guilty of misconduct in carrying on the business}{A business of the firm cannot be carried on except at a loss}{All of the above}
\mcq{494}{It is prepared to find out the result of dissolution, i.e profit or loss:}{Realization account}{Revaluation account}{Profit and loss account}{Profit and loss appropriation account}
\mcq{495}{In case of dissolution, when all the partners are solvent, closing of various assets are credited to:}{Realization account}{Revaluation account}{Profit and loss account}{Assets account}
\mcq{496}{In case of dissolution closing of various liabilities are credited to:}{Realization account}{Revaluation account}{Profit and loss account}{Cash account}
\mcq{497}{In case of dissolution, assets sold for cash are debited to:}{Realization account}{Assets account}{Profit and loss account}{Cash account}
\mcq{498}{In case of dissolution, the payment of liabilities should be credited to:}{Realization}{Cash account}{Profit and loss account}{Liabilities account}
\mcq{499}{In case of dissolution, assets taken over by a partner should be debited to:}{Realization account}{Partner’s capital account}{Asset account}{Cash account}
\mcq{500}{In case of dissolution, the sale of unrecorded asset is debited to:}{Realization account}{Sale account}{Asset account}{Cash account}
\mcq{501}{In case of dissolution, the payment of unrecorded liability is debited to:}{Realization account}{Liability account}{Revaluation account}{Cash account}
\mcq{502}{In case of dissolution, the transfer of undivided profit should be debited to:}{Realization account}{Revaluation account}{Profit and loss account}{Partner’s capital account}
\mcq{503}{The journal entry, for closing of various assets in case of dissolution will be:}{nullRealization a/c..Dr To assets a/c}{nullAsset a/c ….Dr To cash a/c}{nullCash a/c … Dr To asset a/c}{nullCash a/c ..Dr To realization a/c}
\mcq{504}{The journal entry, for closing of various liabilities, in case of dissolution will be:}{nullRealization a/c ..Dr To liabilities a/c}{nullCash a/c … Dr To liabilities a/c}{nullLiabilities a/c …Dr To profit and loss a/c}{nullLiabilities a/c …Dr To realization a/c}
\mcq{505}{The journal entry, for sale of various assets for cash, in case of dissolution will be:}{nullLiabilities a/c …Dr To cash a/c}{nullCash a/c ..Dr To realization a/c}{nullCash a/c … Dr To asset a/c}{nullAsset a/c … Dr To realization a/c}
\mcq{506}{The journal entry, for the payment of liabilities, in case of dissolution will be:}{nullLiabilities a/c …Dr To realization a/c}{nullRealization a/c …Dr To liabilities a/c}{nullCash a/c … Dr To realization a/c}{nullRealization a/c …Dr To cash a/c}
\mcq{507}{The journal entry, for assets taken over by a partner, in case of dissolution will be:}{nullPartner’s capital a/c …Dr To realization a/c}{nullRealization a/c …Dr To liabilities a/c}{nullCash a/c ..Dr To realization a/c}{nullRealization a/c…Dr To assets a/c}
\mcq{508}{The journal entry, for the payment of unrecorded liability, in case of dissolution will be:}{nullLiabilities a/c…Dr To cash a/c}{nullRealization a/c….Dr To cash a/c}{nullCash a/c … Dr To realization a/c}{nullRealization a/c …Dr To liabilities a/c}
\mcq{509}{The journal entry, for the sale of unrecorded asset, in case of dissolution will be:}{nullCash a/c … Dr To asset a/c}{nullAsset a/c… Dr To realization a/c}{nullCash a/c…Dr To realization a/c}{nullRealization a/c…Dr To asset a/c}
\mcq{510}{The journal entry, for a liabilities taken over by a partner, in case of dissolution will be:}{nullRealization a/c…DrTo partner’s capital a/c}{nullLiability a/c …DrTo partner’s capital a/c}{nullLiability a/c… DrTo realization a/c}{nullRealization a/c…DrTo liability a/c}
\mcq{511}{The journal entry, for the transfer of undivided profits, in case of dissolution will be:}{nullRealization a/c..DrTo profit a/c}{nullPartner’s capital a/c…DrTo realization a/c}{nullRealization a/c…DrTo partner’s capital a/c}{nullProfit and loss a/c …DrTo partner’s capital a/c}
\mcq{512}{The journal entry, for expense on realization will be:}{nullRealization a/c …DrTo cash a/c}{nullExpenses a/c…DrTo cash a/c}{nullExpenses a/c…DrTo profit \& loss a/c}{nullProfit \& loss a/c..DrTo expenses a/c}
\mcq{513}{The journal entry, for transferring of loss on realization will be:}{nullPartner’s capital a/c…DrTo profit \& loss a/c}{nullPartner’s capital a/c…DrTo realization a/c}{nullRealization a/c…DrTo partner’s capital a/c}{nullProfit \& loss a/c…DrTo realization a/c}
\mcq{514}{The journal entry, for repayment of partner’s capital balance sheet in case dissolution will be:}{nullCash a/c… DrTo partner’s capital a/c}{nullPartner’s capital a/c…DrTo cash a/c}{nullRealization a/c…DrTo partner’s capital a/c}{nullProfit \& loss a/c …DrTo realization a/c}
\mcq{515}{The journal entry, for the amount brought in by partner, in case of dissolution will be:}{nullCash a/c… DrTo partner’s capital a/c}{nullPartner’s capital a/c…DrTo cash a/c}{nullRealization a/c…DrTo partner’s capital a/c}{nullPartner’s capital a/c…DrTo realization a/c}
\mcq{516}{The entry for transferring partner’s capital to loan a/c will be:}{nullCash a/c… DrTo partner’s a/c}{nullCash a/c … DrTo realization a/c}{nullPartner’s capital a/c…DrTo partner’s loan a/c}{nullPartner’s loan a/c ..DrTo realization a/c}
\mcq{517}{The journal entry for the repayment of partner’s loan in case of dissolution will be:}{nullPartner’s loan a/c…DrTo cash a/c}{nullRealization a/c …DrTo partner’s loan a/c}{nullPartner’s loan a/c…DrTo partner’s capital a/c}{nullNo entry}
\mcq{518}{Reserve for bad and doubtful appearing in the books at the time of dissolution is transferred to:}{Debtors account}{Revaluation account}{Reserve for bad debts account}{Realization account}
\mcq{519}{General reserve or accumulated profits at the time of dissolution is transferred to:}{Realization account}{Solvent partners’ capital account}{Insolvent partner’s capital account}{All partners’ capital accounts}
\mcq{520}{Joint life policy reserve at the time of dissolution is transferred to:}{Realization account}{Solvent partners’ capital}{Insolvent partner’s capital account}{All partner’s capital account}
\mcq{521}{According to the decision of Garner Vs Murray rule, the loss of insolvent partner is to be shared by solvent partners in the:}{New profit sharing ratio}{Old profit sharing ratio}{Capital ratio}{Gaining ratio}
\mcq{522}{In case of dissolution, partners personal assets are first used for the payment of:}{Personal liabilities}{Firm’s liabilities}{Court liabilities}{None of these}
\mcq{523}{The balance left in the capital accounts in case of dissolution is transferred to:}{Revaluation account}{Realization account}{Bank account}{Profit \& loss appropriation account}
\mcq{524}{In case of dissolution, the loss due to insolvency of partner, wen capitals are fixed is to be shared by the solvent partners in:}{Profit sharing ratio}{Fixed capital sharing ratio}{Drawing ratio}{Current account ratio}
\mcq{525}{Due to the insolvency of all partner’s the journal entry to close the capital account against the profit and loss account will be:}{nullProfit \& loss a/c…DrTo partner’s capital a/c}{nullProfit \& loss a/c…DrTo realization a/c}{nullRealization a/c…DrTo profit \& loss a/c}{nullProfit \& loss a/c…DrTo cash a/c}
\end{document}
